\usepackage{fontspec}
\IfFileExists{UEOT_CJK.ttf}{
  \setmainfont{UEOT_CJK.ttf}[BoldFont=UEOT_CJK.ttf,BoldFeatures={FakeBold=2},ItalicFont=UEOT_CJK.ttf,ItalicFeatures={FakeSlant=.15}]
  \setsansfont{UEOT_CJK.ttf}[BoldFont=UEOT_CJK.ttf,BoldFeatures={FakeBold=2}]
  \setmonofont{UEOT_CJK.ttf}
}{
  \setmainfont{Noto Sans CJK SC}
  \setsansfont{Noto Sans CJK SC}
  \setmonofont{Noto Sans CJK SC}
}
\XeTeXlinebreaklocale "zh"
\XeTeXlinebreakskip=0pt plus 1pt
\usepackage{mathtools}
\usepackage{fvextra}
\usepackage{etoolbox}
\usepackage{newunicodechar}
\newunicodechar{ℏ}{\ensuremath{\hbar}}
\newunicodechar{∎}{\ensuremath{\square}}
\usepackage{fancyhdr}
\pagestyle{fancy}
\fancyhf{}
\fancyhead[L]{\small UEOT PHYSICS}
\fancyhead[R]{\small 物理学 · 第三版}
\fancyfoot[C]{\small\thepage}
\setlength{\headheight}{16pt}
\renewcommand{\headrulewidth}{0.3pt}
\setcounter{tocdepth}{0}
\setlength{\parindent}{0pt}
\setlength{\parskip}{5pt plus 1pt}
\setlength{\emergencystretch}{8em}
\setlength{\tabcolsep}{3pt}
\renewcommand{\arraystretch}{1.16}
\AtBeginEnvironment{longtable}{\small}
\allowdisplaybreaks[2]
\renewcommand{\contentsname}{目录}
\makeatletter
\renewcommand{\maketitle}{
  \begin{titlepage}
  \thispagestyle{empty}
  \vspace*{30mm}
  {\color{teal}\fontsize{46}{52}\selectfont UEOT\par}
  \vspace{5mm}
  {\fontsize{32}{40}\selectfont\bfseries 物理学\par}
  \vspace{11mm}
  {\Large 从对象数学到量子物理\par}
  \vspace{4mm}
  {\large 理论推导、形式证明与研究边界\par}
  \vspace{10mm}
  {\color{teal}\rule{\textwidth}{1pt}\par}
  \vspace{7mm}
  {\large 第三版\quad 3.0 正式整合版\par}
  \vspace{6mm}
  {钱德桂（Degui Qian）\par}
  \vspace{5mm}
  {51章\quad 6个附录\quad UEOT-QM Kernel 1.4\par}
  \vfill
  {对象与协议\quad 量子操作与场论\quad 证明与反例\par}
  \vspace{7mm}
  {2026年10月20日\par}
  \end{titlepage}
}
\makeatother
